% GENERATED FILE. Edit canonical lessons, then run npm run generate:books.
% canonical-source-hash: fnv1a64:4c177339d3e76ff3
% canonical-lessons: AR-C189-muhasib, AR-C189-suhufi, AR-C189-amil, AR-C189-mutarjim, AR-C189-muqabala

\chapter{Accountant, Journalist, Worker, Translator, Interview}
\label{ch:ar-accountant-journalist-worker-translator-interview}
\hlchaptermodality{189}

\begin{chapteropening}
\textbf{By the end of this chapter:} \emph{I can say an accountant, a journalist, a worker, a translator and an interview.}

Complete the last lesson of chapter 189: I can say an accountant, a journalist, a worker, a translator and an interview.
\end{chapteropening}

\section[\texorpdfstring{\ar{محاسب} (muḥāsib)}{محاسب (muḥāsib)}]{\ar{محاسب} (muḥāsib) — an accountant}

\label{lesson:AR-C189-muhasib}



\begin{quote}
Before the new one: say the Arabic for a customer, then the Arabic for a cashier; a money changer.
\end{quote}

\subsection*{You'll want to know: \ar{محاسب}}
\textbf{\ar{محاسب}} — \emph{muḥāsib} — \textquotedblleft{}an accountant\textquotedblright{}.

\begin{cousinweb}
\textbf{\ar{محاسب}} (\emph{muḥāsib}) is built on the native Arabic root \emph{ḥ-s-b}, \textquotedblleft{}to count, to reckon\textquotedblright{} (Lane). \emph{Muḥāsib} is the Form III participle, \textquotedblleft{}one who settles accounts\textquotedblright{}.
\end{cousinweb}

\begin{grammarlens}[title={What you've built}]
One more word for this chapter.
\end{grammarlens}

\subsection*{Your turn}
\begin{itemize}
\raggedright
  \item \emph{Say it:} \emph{muḥāsib}
  \item \emph{Say it:} \emph{muḥāsib}, once more
  \item \emph{Say it:} say \emph{muḥāsib}
  \item \emph{Recall:} say the Arabic for a customer, then the Arabic for a cashier; a money changer, then say \emph{muḥāsib} again
\end{itemize}

\begin{tcolorbox}[breakable,colback=teal!4,colframe=teal!35!black,title={Before you move on}]
What does \ar{محاسب} mean? (\textquotedblleft{}An accountant\textquotedblright{}.) Say it once more.
\end{tcolorbox}
\section[\texorpdfstring{\ar{صحفي} (ṣuḥufī)}{صحفي (ṣuḥufī)}]{\ar{صحفي} (ṣuḥufī) — a journalist}

\label{lesson:AR-C189-suhufi}



\begin{quote}
Before the new one: say the Arabic for a cashier; a money changer, then the Arabic for an accountant.
\end{quote}

\subsection*{You'll want to know: \ar{صحفي}}
\textbf{\ar{صحفي}} — \emph{ṣuḥufī} — \textquotedblleft{}a journalist\textquotedblright{}.

\begin{cousinweb}
\textbf{\ar{صحفي}} (\emph{ṣuḥufī}) is built on \emph{ṣaḥīfa}, \textquotedblleft{}a sheet, a page\textquotedblright{}, a native noun, filed under the root \emph{ṣ-ḥ-f} (Lane). \emph{Ṣuḥufī} is \emph{ṣuḥuf}, \textquotedblleft{}sheets, papers\textquotedblright{}, with the adjective ending \emph{-ī}; it is a modern word.
\end{cousinweb}

\begin{grammarlens}[title={What you've built}]
One more word for this chapter.
\end{grammarlens}

\subsection*{Your turn}
\begin{itemize}
\raggedright
  \item \emph{Say it:} \emph{ṣuḥufī}
  \item \emph{Say it:} \emph{ṣuḥufī}, once more
  \item \emph{Say it:} say \emph{ṣuḥufī}
  \item \emph{Recall:} say the Arabic for a cashier; a money changer, then the Arabic for an accountant, then say \emph{ṣuḥufī} again
\end{itemize}

\begin{tcolorbox}[breakable,colback=teal!4,colframe=teal!35!black,title={Before you move on}]
What does \ar{صحفي} mean? (\textquotedblleft{}A journalist\textquotedblright{}.) Say it once more.
\end{tcolorbox}
\section[\texorpdfstring{\ar{عامل} (ʿāmil)}{عامل (ʿāmil)}]{\ar{عامل} (ʿāmil) — a worker}

\label{lesson:AR-C189-amil}



\begin{quote}
Before the new one: say the Arabic for an accountant, then the Arabic for a journalist.
\end{quote}

\subsection*{You'll want to know: \ar{عامل}}
\textbf{\ar{عامل}} — \emph{ʿāmil} — \textquotedblleft{}a worker\textquotedblright{}.

\begin{cousinweb}
\textbf{\ar{عامل}} (\emph{ʿāmil}) is built on the native Arabic root \emph{ʿ-m-l}, \textquotedblleft{}to work\textquotedblright{} (Lane). \emph{ʿĀmil} is the active participle, \textquotedblleft{}one who works\textquotedblright{}.
\end{cousinweb}

\begin{grammarlens}[title={What you've built}]
One more word for this chapter.
\end{grammarlens}

\subsection*{Your turn}
\begin{itemize}
\raggedright
  \item \emph{Say it:} \emph{ʿāmil}
  \item \emph{Say it:} \emph{ʿāmil}, once more
  \item \emph{Say it:} say \emph{ʿāmil}
  \item \emph{Recall:} say the Arabic for an accountant, then the Arabic for a journalist, then say \emph{ʿāmil} again
\end{itemize}

\begin{tcolorbox}[breakable,colback=teal!4,colframe=teal!35!black,title={Before you move on}]
What does \ar{عامل} mean? (\textquotedblleft{}A worker\textquotedblright{}.) Say it once more.
\end{tcolorbox}
\section[\texorpdfstring{\ar{مترجم} (mutarjim)}{مترجم (mutarjim)}]{\ar{مترجم} (mutarjim) — a translator}

\label{lesson:AR-C189-mutarjim}



\begin{quote}
Before the new one: say the Arabic for a journalist, then the Arabic for a worker.
\end{quote}

\subsection*{You'll want to know: \ar{مترجم}}
\textbf{\ar{مترجم}} — \emph{mutarjim} — \textquotedblleft{}a translator\textquotedblright{}.

\begin{cousinweb}
\textbf{\ar{مترجم}} (\emph{mutarjim}) is the participle of \emph{tarjama}, \textquotedblleft{}to translate\textquotedblright{}, which Lane gives with \emph{tarjumān}, \textquotedblleft{}interpreter\textquotedblright{}. \emph{Tarjumān} is usually traced to Aramaic \emph{targmānā}; English \emph{dragoman} comes from it.
\end{cousinweb}

\begin{grammarlens}[title={What you've built}]
One more word for this chapter.
\end{grammarlens}

\subsection*{Your turn}
\begin{itemize}
\raggedright
  \item \emph{Say it:} \emph{mutarjim}
  \item \emph{Say it:} \emph{mutarjim}, once more
  \item \emph{Say it:} say \emph{mutarjim}
  \item \emph{Recall:} say the Arabic for a journalist, then the Arabic for a worker, then say \emph{mutarjim} again
\end{itemize}

\begin{tcolorbox}[breakable,colback=teal!4,colframe=teal!35!black,title={Before you move on}]
What does \ar{مترجم} mean? (\textquotedblleft{}A translator\textquotedblright{}.) Say it once more.
\end{tcolorbox}
\section[\texorpdfstring{\ar{مقابلة} (muqābala)}{مقابلة (muqābala)}]{\ar{مقابلة} (muqābala) — an interview}

\label{lesson:AR-C189-muqabala}



\begin{quote}
Before the new one: say the Arabic for a worker, then the Arabic for a translator.
\end{quote}

\subsection*{You'll want to know: \ar{مقابلة}}
\textbf{\ar{مقابلة}} — \emph{muqābala} — \textquotedblleft{}an interview\textquotedblright{}.

\begin{cousinweb}
\textbf{\ar{مقابلة}} (\emph{muqābala}) is built on the native Arabic root \emph{q-b-l}, \textquotedblleft{}to face\textquotedblright{} (Lane). \emph{Muqābala} is the verbal noun of Form III, \textquotedblleft{}meeting face to face\textquotedblright{}.
\end{cousinweb}

\begin{grammarlens}[title={What you've built}]
One more word for this chapter.
\end{grammarlens}

\subsection*{Your turn}
\begin{itemize}
\raggedright
  \item \emph{Say it:} \emph{muqābala}
  \item \emph{Say it:} \emph{muqābala}, once more
  \item \emph{Say it:} say \emph{muqābala}
  \item \emph{Recall:} say the Arabic for a worker, then the Arabic for a translator, then say \emph{muqābala} again
\end{itemize}

\begin{tcolorbox}[breakable,colback=teal!4,colframe=teal!35!black,title={Before you move on}]
What does \ar{مقابلة} mean? (\textquotedblleft{}An interview\textquotedblright{}.) Say it once more.
\end{tcolorbox}
